\section{模型评价、改进与推广}

\subsection{模型优点}
本文四问覆盖从数据刻画到能力预测的主要环节，但数值链条仍有断点：问题一的域级质量用于问题三的质量成本基线；配比贡献尚未进入广义律数值计算；损失到能力的桥接是单独拟合的探索性映射。A1 固定参照与逐条复核使质量评分可重算；广义律在 $Q=1$ 时满足代数退化。优化部分用 KKT 刻画结构、预算消元与质量剖面搜索候选解，并以份额导数峰值定义候选结构转移。对半合成质量标度方向异常和前沿窗口规模负增量均保留并讨论。

\subsection{模型不足与改进方向}
其一，质量放大指数无法从唯一一组半合成数据识别，其经验拟合落在边界且口径与部署质量不同，故下游采用了反映“质量有益”前提的取值并做敏感性；若能获得大规模真实的质量--损失联合观测，可把该指数从假设升级为估计。其二，配比没有进入当前标度律数值计算，也未联立优化 $17$ 维配比；改进可把配比加入决策变量并以投影梯度处理单纯形约束，代价是转移分析的可解释性下降。其三，前沿预测依赖能力斜率情景假设，并未估计算力增速到能力斜率的映射，改进可引入更细的算力供给模型或以真实新发布模型滚动校正分位带。

\subsection{推广}
广义标度律的“有效数据量”形式可推广到其它可归一化的数据属性（如去重率、领域纯度），并在归一化边界 $Q=1$ 退化为经典形式。问题三的三部分算力约束可扩展到含推理成本或多阶段训练的预算结构。问题四的描述性分解与分层桥接可用于其它基准体系；这些推广均需在目标场景重新标定参数。

\section{结论}
本文用语义解码、A1 固定百分位参照和三组平衡权重重算问题一，得七域质量中位数 $Q_0=0.48788$；按高教育价值且判有广告的规则，A1 冲突率为 $0.303\%$，其中 C4 的域内冲突率最高，为 $1.03\%$。配比模型在同尺度检验上具有一定解释力，但跨尺度绝对损失预测失准。经典标度律拟合良好；广义律的质量指数仍属结构性假设。在新 $Q_0$ 下，指数型质量成本的三档预算候选最低损失为 $3.073$、$2.151$、$1.914$，质量份额变化峰值约在 $10^{19}$ FLOPs；注意力与基础训练开销相当的上下文长度为 $3\times10^{4}$ tokens。能力前沿预测依赖情景斜率，不能视为已校准的无条件预测。质量分数与资源配置结论的主要限制是评分权重及成本函数尚无训练损失标签的联合校准。
